\section{Introduction}
\label{sec:intro}

LLM scientific ideation has crossed a surprising threshold. In a controlled study with more than 100 NLP researchers, ideas produced by an LLM agent were judged \emph{more novel} than those written by expert humans~\citep{si2024can}. Yet the field still distrusts machine-generated ideas, and for a concrete reason: ``novelty'' is ill-defined. Unguided LLMs swing between incremental rephrasings of known work and implausible over-reaches, and the very judgments that crown them as novel come from other LLMs whose reliability is in question.

A natural and inexpensive fix is to \emph{ground} the abstract notion of novelty in a concrete reasoning process before an idea is committed. Rather than asking a model to ``be novel,'' one can ask it to flip a stated assumption, induce a hypothesis from observations, draft a full proposal, imagine the most surprising plausible outcome, design a falsification test, or explicitly contrast its idea against retrieved prior work. Each of these is a cheap prompting protocol, not a new model, so if any reliably moves ideas toward the novel-and-plausible sweet spot it would directly improve automated discovery pipelines at almost no cost.

\para{The gap.} Every strong prior system champions \emph{one} grounding mechanism, evaluated on \emph{its own} dataset and metric: literature comparison in SciMON, Chain-of-Ideas, and Nova~\citep{wang2023scimon,li2024chain,nova2024}; induction in MOOSE and causal-graph hypothesis generation~\citep{yang2023moose,causalgraph2024}; deduction in HypoGen's Bit$\rightarrow$Flip$\rightarrow$Spark schema~\citep{hypogen2025}; proposal writing in the Si agent and VirSci~\citep{si2024can,virsci2024}; outcome imagination in AutoDiscovery's Bayesian surprise~\citep{autodiscovery2025}; and critic/small-experiment feedback in Deep-Ideation~\citep{deepideation2025}. \emph{No paper compares these six mechanisms head-to-head under one controlled protocol}---same backbone, same seed topics, same metrics, same output format. Two evaluation-skeptic papers~\citep{glitters2025,rqbench2026} add that the \emph{evaluation} itself must be grounded, not left to a single LLM's opinion. Whether grounding helps is therefore essentially untested as a controlled ablation, and how to measure ``help'' is itself contested.

\para{Our approach.} We run the first controlled head-to-head ablation of all six grounding mechanisms plus an unguided control. We hold the backbone (\gptfourone), the seed topics (the 7 canonical NLP topics of \citet{si2024can}), and the output schema fixed, varying \emph{only} the reasoning scaffold that precedes the idea; this isolates grounding from the length and format confounds that complicate prior comparisons. We generate 5 replicates per cell, for 245 ideas, and evaluate them through three independent lenses: (i)~an \emph{objective} literature-grounded novelty metric (\lgn), the embedding distance from each idea to the nearest of 1{,}050 real prior-work abstracts retrieved from Semantic Scholar; (ii)~a 2-judge rubric using \claudesonnet and \geminiflash, both deliberately off-family from the backbone to limit self-preference; and (iii)~position-controlled pairwise comparisons against the control.

\para{What we find.} The three lenses disagree, and the disagreement is the headline. Grounding does not uniformly improve quality---it \emph{steers} the novelty--feasibility trade-off. Deduction and literature comparison raise judged novelty (Cohen's $d = 0.81$ and $0.57$), and literature comparison wins 79\% of head-to-head comparisons against the control while being the only mechanism that improves idea-set diversity. But \emph{no} mechanism increases objective distance from prior work; the only significant effect there is proposal writing making ideas \emph{more} conventional. Most strikingly, LLM-judged novelty and objective novelty are uncorrelated at the item level (Spearman $\rho = 0.01$, $p = 0.85$), even though the two judges agree with each other ($\rho = 0.41$). The novelty gains are real in the judge's eye and a mirage under a literature-anchored measure---a clean, controlled reproduction of the ``novelty mirage''~\citep{rqbench2026}.

\para{Contributions.}
\begin{itemize}[leftmargin=*,itemsep=0pt,topsep=0pt]
    \item We conduct the first controlled head-to-head ablation of six grounding mechanisms for LLM idea generation, isolating the reasoning scaffold from backbone, topic, and output-format confounds across 245 ideas.
    \item We complement the standard LLM-judge rubric with an objective, judge-independent literature-grounded novelty metric and a convergent-validity test, and show that judged novelty is uncorrelated with objective novelty---a direct, quantitative confirmation of the novelty mirage.
    \item We characterize each mechanism's effect as a position on the novelty--feasibility axis, and distill practical guidance: use literature comparison for judged novelty and diversity, induction or proposal writing for feasibility, and never report single-LLM-judge novelty alone.
\end{itemize}

The rest of the paper reviews related work (\secref{sec:related}), details our protocol and metrics (\secref{sec:method}), presents results across the three lenses (\secref{sec:results}), and discusses implications and limitations (\secref{sec:discussion}).
